% !TEX program = lualatex
% ECON 201, Worksheet 7: Demand and Elasticity. Compile with LuaLaTeX so the
% PDF is tagged (see syllabus/Econ201-Syllabus.tex for the same setup).
% Build from the course root:
%   latexmk -lualatex -cd worksheets/worksheet07.tex && latexmk -c -cd worksheets/worksheet07.tex
% Every number here is computed and printed by slides/figures/make_lecture7.py
% so the worksheet and the deck quote the same figures (the class Dodgers demand).
\DocumentMetadata{
  pdfstandard=UA-2,
  pdfversion=2.0,
  lang=en-US,
  tagging=on
}
\documentclass[11pt]{article}

\usepackage[letterpaper, margin=0.8in, top=0.6in, bottom=0.55in]{geometry}
\usepackage{fontspec}
\setmainfont{Fira Sans}
\usepackage{amsmath}
\usepackage{amssymb}
\usepackage{graphicx}
\usepackage{booktabs}
\usepackage{array}
\usepackage{xcolor}
\usepackage{needspace}
\usepackage[most]{tcolorbox}
\definecolor{cream}{HTML}{FDF3EA}
\definecolor{accent}{HTML}{BF5700}
\usepackage{titlesec}
\titleformat{\section}{\large\bfseries\color{accent}}{}{0pt}{}
\titlespacing*{\section}{0pt}{6pt}{2pt}
\usepackage{hyperref}
\hypersetup{pdftitle={ECON 201 Worksheet: Demand and Elasticity}, pdfauthor={Div Bhagia}, hidelinks}
\pagestyle{empty}
\setlength{\parindent}{0pt}
\setlength{\parskip}{3pt}

% Ruled answer space: n lines. Plain TeX loop; pgffor's \foreach breaks the
% enumerate counter when used inside an item.
\newcount\answerlinecount
\newcommand{\answerlines}[1]{%
  \par\vspace{1pt}\answerlinecount=#1
  \loop\ifnum\answerlinecount>0
    \rule{\linewidth}{0.3pt}\par\vspace{5pt}%
    \advance\answerlinecount by -1
  \repeat
  \vspace{-5pt}}

% Ruled grids rather than booktabs tables: students write in the cells, so
% every cell gets visible borders and the blank rows get extra height. One
% column layout for all three activities, so the tables line up.
\newcolumntype{C}{>{\centering\arraybackslash}p{3.0cm}}
\newcommand{\blankrow}{\rule{0pt}{1.8em}}

\begin{document}

{\LARGE\bfseries\color{accent} Worksheet: Demand and Elasticity}\hfill{\large ECON 201}

% Terms sit on the cream panel, where the brand orange falls below 4.5:1;
% the darker orange keeps WCAG AA contrast (see CLAUDE.md).
\definecolor{accentdark}{HTML}{8F4100}
\newcommand{\term}[1]{\textcolor{accentdark}{\textbf{#1}}}

\begin{tcolorbox}[colback=cream, colframe=accent, boxrule=0.6pt, arc=2pt,
  left=6pt, right=6pt, top=4pt, bottom=4pt, title=Definitions,
  fonttitle=\bfseries, coltitle=white]
\term{Percentage change}:
\[
  \%\text{ change} = \frac{\text{new} - \text{old}}{\text{old}} \times 100
\]
\term{Price elasticity of demand}, $\varepsilon$: the percentage change in
demand that would occur in response to a 1\% increase in price,
\[
  \varepsilon = -\frac{\%\text{ change in } Q}{\%\text{ change in } P}
\]
If price changes by $\Delta P$ and demand changes by $\Delta Q$, three other
equivalent ways to write it are
\[
  \varepsilon = -\frac{\Delta Q / Q}{\Delta P / P}
  = -\frac{P}{Q} \times \frac{\Delta Q}{\Delta P}
  = -\frac{P}{Q} \times \frac{1}{\text{slope}}
\]
where the slope of the demand curve is $\Delta P / \Delta Q$.

Demand is \term{elastic} if $\varepsilon$ is greater than 1, and
\term{inelastic} if $\varepsilon$ is less than 1.
\end{tcolorbox}

In Lecture 3, you told us the most you would pay for a day at a Dodgers game.
Linearizing your answers gives this demand curve:
\[
  Q = 90 - 0.3\,P
\]

\needspace{14\baselineskip}\section{1. Elasticity from percentage changes}

Use percentage changes to find the elasticity of demand for each \$10 price
rise below.

\begin{center}
\renewcommand{\arraystretch}{1.5}
\newcolumntype{D}{>{\centering\arraybackslash}p{2.7cm}}
\begin{tabular}{|>{\raggedright\arraybackslash}p{2.9cm}|D|D|D|D|}
\hline
$P$ rises from & $Q$ falls from & \% change in $P$ & \% change in $Q$ & Elasticity, $\varepsilon$ \\
\hline
\$50 to \$60 & \rule{0pt}{1.8em} & & & \\
\hline
\$100 to \$110 & \rule{0pt}{1.8em} & & & \\
\hline
\$200 to \$210 & \rule{0pt}{1.8em} & & & \\
\hline
\end{tabular}
\end{center}

At which prices is demand elastic?
\answerlines{1}

\needspace{10\baselineskip}\section{2. Elasticity from the formula}

Use $\varepsilon = -\dfrac{P}{Q} \times \dfrac{\Delta Q}{\Delta P}$ to find the
elasticity at the same three prices. With $Q = 90 {\color{accent}\,-\,0.3}\,P$,
$\Delta Q / \Delta P = -0.3$.

\begin{center}
\renewcommand{\arraystretch}{1.5}
\begin{tabular}{|>{\raggedright\arraybackslash}p{4.2cm}|C|C|C|}
\hline
Price, $P$ & \$50 & \$100 & \$200 \\
\hline
Demand, $Q$ & 75 & 60 & 30 \\
\hline
Elasticity, $\varepsilon$ & \rule{0pt}{1.8em} & & \\
\hline
\end{tabular}
\end{center}

Do you get the same answers as in Activity 1?
\answerlines{1}

\end{document}
